\subsection{Stun systém a dočasné zpomalení}

Když projektil zasáhne kouli, volá se metoda \texttt{Stun()}. Ta inkrementuje counter \newline \texttt{\_activeStuns} a spustí korutinu \texttt{StunCooldown()}. Systém rozlišuje první stun a další stuny -- u prvního se uloží aktuální rychlost jako základ, u dalších se jen nastaví nízká rychlost.

Více stunů současně způsobuje, že rychlost zůstává na minimální hodnotě, dokud se všechny nestihnou vypořádat. Po uplynutí \texttt{stunDuration} se counter sníží a rychlost se obnoví na základní úroveň. Pokud je více stunů aktivních, rychlost zůstává na minimu.

Tento systém umožňuje hráči střílet více projektilů za sebou pro delší zpomalení, ale vzhledem k omezené munici (2 náboje) je toto omezeno. Navíc po každém stunu se základní rychlost zvýší, takže koule je postupně rychlejší.

\begin{lstlisting}[language=csharp,caption={Listing 15: Metoda Stun a StunCooldown}]
public void Stun(bool stun)
{
    if (!stun) return;
    _activeStuns++;
    StartCoroutine(StunCooldown());
}

private IEnumerator StunCooldown()
{
    if (_activeStuns == 1) _baseMoveSpeed = moveSpeed;
    moveSpeed = 5f;

    yield return new WaitForSeconds(stunDuration);

    _activeStuns = Mathf.Max(0, _activeStuns - 1);
    _baseMoveSpeed = Mathf.Min(_baseMoveSpeed + speedIncreasePerStun, maxSpeed);
    moveSpeed = _activeStuns == 0 ? _baseMoveSpeed : 5f;
}
\end{lstlisting}
